% Complete local proofs for the explicitly declared quotient presentation.
\section{Imaginaries and quotient interpretations}

\begin{annexitem}{annex:p03-imaginaries-quotient-truth}{thm:p03-imaginaries-quotient-truth}{Descent, term values, and every satisfaction clause}
\begin{proof}
An equivalence relation on the nonempty set \(D_s(M)\) has a nonempty
quotient set. For a relation formula, representative invariance gives exactly
one truth value on each tuple of classes. For a function graph, take any
representatives of its finitely many input classes. An output exists by the
totality condition. Any two outputs are equivalent, and changing any input or
output representative leaves the graph true. Thus the resulting output class
exists and is independent of every choice. This defines a function on the
classes. A constant is the zero-input case. Nonemptiness and these assertions
include zero-dimensional sorts: there is one empty representative tuple.

Induct on a term \(t\). For a variable, \(V_t^I\) says precisely that
its output representative is equivalent to the assigned representative.
For a constant this is its invariant graph and unique output class. In the
compound case the induction hypothesis makes each existentially chosen
\(\bar u_i\) a representative of the actual value of \(t_i\).
The function graph then makes \(\bar z\) a representative of the value
of \(f(t_1,\ldots,t_k)\). Conversely choose representatives for those
finitely many term values and apply totality. This is a finite sequence of
existence choices, not an application of the Axiom of Choice to an infinite
family. The compound formula therefore selects all and only the representatives
of its term value.

For an atomic relation, this term induction and invariance give exactly its
truth in the interpreted structure. For equality, two term values are the
same quotient element exactly when their representatives satisfy \(E_s\).
A nullary relation uses its assigned sentence, and falsity is false on both
sides. Negation, conjunction and disjunction follow from their truth clauses
and the induction hypothesis; implication follows in the same way from its
Boolean clause. For an existentially quantified variable, a witness class has
a representative in \(D_s\), and a representative satisfying the translated
matrix supplies a witness class. The guard is thus exactly the target range.
For a universal, every representative gives a class, every class has a
representative, and invariance makes the translated matrix independent of
which representative is used. This proves both directions of the universal
clause. A zero-coordinate quantified block is the same argument on the
singleton quotient, with no source variable to quantify.
Structural induction proves the truth lemma for every formula and every domain
assignment. Applying it to the translated axioms proves \(I(M)\models U\).
\end{proof}
\end{annexitem}
\begin{annexitem}{annex:p03-imaginaries-complexification}{thm:p03-imaginaries-complexification}{Coordinate operations and the nonzero norm}
\begin{proof}
Coordinate addition is an abelian group operation, with the displayed zero
and negation. The multiplication formula is symmetric, so commutative.
Multiplying three pairs \((a,b),(c,d),(e,f)\) in either association gives
real coordinate \(ace-bde-adf-bcf\) and imaginary coordinate
\(acf-bdf+ade+bce\), by the field distributive laws. Thus it is
associative, and \((1,0)\) is the identity. Expanding the two coordinates
also gives distributivity over coordinate addition.
In an ordered field, a nonzero square is positive and a square is nonnegative.
For \((a,b)\ne(0,0)\), at least one of \(a^2,b^2\) is positive,
so \(r=a^2+b^2>0\). Multiplication by \((a/r,-b/r)\) gives
\((r/r,0)=(1,0)\). This proves every field axiom, with distinct zero and
one since \(F\) is a field. The inverse graph is the first-order condition
\(r\ne0\land ru=a\land rv=-b\), with \(r=a^2+b^2\);
existence and uniqueness use division in \(F\), not an extra primitive.
The map \(a\mapsto(a,0)\) preserves both operations and constants and is
injective. Multiplication makes \((0,1)^2=(-1,0)\).
For the real field the identification \((a,b)\mapsto a+bi\) preserves
exactly these coordinate operations, giving the usual complex field.
The proof uses neither a polynomial-root theorem nor a claim of algebraic
closure, and therefore does not silently turn arbitrary ordered fields into
real-closed fields.
\end{proof}
\end{annexitem}

\begin{annexitem}{annex:p03-imaginaries-elementarity}{thm:p03-imaginaries-elementarity}{The natural quotient maps, not literal class inclusion}
\begin{proof}
Elementarity of \(h\) preserves and reflects each \(D_s\) and \(E_s\).
Equivalent representatives therefore have equivalent images; conversely an
equivalence of images implies equivalence of the originals. The formula for
\(h^I\) is consequently well defined and injective. Each relation formula
is preserved and reflected. A true function graph on representatives remains
true after applying \(h\), so the quotient function commutes with \(h^I\);
the same observation handles constants. For an arbitrary \(K\)-formula,
the truth lemma translates it into an \(L\)-formula. Elementarity of
\(h\), followed by the truth lemma in \(N\), proves elementarity of
\(h^I\).

The elementary-chain union is elementary over each \(M_i\): the usual
formula induction has its existential step at a stage containing both the
finite parameter tuple and the chosen witness, after which elementarity
reflects the existential sentence to the earlier stage. Atomic interpretations
are compatible unions; Boolean cases are immediate; the universal case is
the negated existential case. This also proves the required union assertion
without an additional unspecified chain lemma.

Form the directed limit sort from pairs \((i,e_i)\), where
\(e_i\in I(M_i)_s\), identifying two pairs if their images agree in a
common later stage. Send the pair represented by \(\bar a\in D_s(M_i)\)
to \([\bar a]_{E_s(M_\infty)}\). The preceding elementary-map result
makes this well defined. It is onto: a finite representative tuple in the union
lies in one stage, and membership in \(D_s\) reflects to that stage.
For injectivity, put two representative tuples in one common stage. If they
are equivalent in the union, elementarity reflects their \(E_s\)-formula
to that stage, making the two limit pairs equivalent. Relations and operation
graphs are transported by the same argument. The map is therefore an
isomorphism. Each canonical map to it is elementary by the first part, applied
to \(M_i\preccurlyeq M_\infty\). No equality between the underlying
sets of successive equivalence classes was assumed.
\end{proof}
\end{annexitem}

\begin{annexitem}{annex:p03-imaginaries-composition}{thm:p03-imaginaries-composition}{Flattening representatives and equality}
\begin{proof}
For a target sort \(s\), let \(P_s\) be the finite product of all
\(J\)-representative domains needed by the coordinates of an
\(I\)-representative. Their quotient projections give an onto map from
\(P_s\) to that coordinate product in \(J(M)\). Pull back \(D_s^I\)
along this map. By the truth lemma for \(J\), the pullback is exactly the
flattened domain stated in the theorem. A nonempty \(D_s^I\) has a
representative there, using only finitely many choices. Two flattened tuples
satisfy \((E_s^I)^J\) precisely when their projected tuples are
\(E_s^I\)-equivalent. This proves reflexivity, symmetry and transitivity
of the flattened equality and gives a bijection of its quotient with
\(I(J(M))_s\).

The same truth lemma makes every flattened relation graph invariant and every
flattened function graph total and single-valued modulo this equality.
The bijections therefore preserve all relations, functions and constants.
For formulas, atomic equality and relation cases have just been established;
the term-value induction introduces exactly the intermediate representative
blocks needed by the two successive graph evaluations. Boolean clauses
commute. An interpreted existential first requires an \(I\)-domain tuple
in \(J(M)\); representing its coordinates by \(J\) requires exactly
the flattened block guards and the translated \(I\)-domain formula.
The universal clause uses those same guards by implication. These facts prove
agreement with iterated translation by structural induction, after consistently
renaming all bound variables. They also prove every translated target axiom.
Different coordinate sorts contribute their respective block lengths, so the
dimension assertion is their finite sum. At no point is a definable section
of either quotient chosen.
\end{proof}
\end{annexitem}

\begin{annexitem}{annex:p03-imaginaries-eq-pullback}{thm:p03-imaginaries-eq-pullback}{All quotient sorts, restricted domains, and parameter fibers}
\begin{proof}
For each imaginary sort use the whole home representative product as domain
and the indicated \(E\) as quotient equality. A home sort uses actual
equality. The graph of \(q_E\) is interpreted by \(E\) between its
home input tuple and its output representative tuple. The declared extra
relations have invariant pullbacks by definition. Thus these are quotient
interpretation data. Each formula mentions finitely many sorts, so the term
and formula proof already given applies to it, proving pullback and elementary
preservation even when the total language has infinitely many sorts.

For \(E^*\), the equivalence classes inside \(D\) are exactly the
old \(E\)-classes, and every tuple outside \(D\) has its own singleton
class. There is no class meeting both regions. This explicitly proves the
equivalence-relation axioms. Its quotient subset is defined in
\(M^{\mathrm{eq}}\) by
\(\exists\bar x(D(\bar x)\land q_{E^*}(\bar x)=z)\).
The quotient map from \(D/E\) to this subset is bijective. This formula
defines the empty subset if \(D\) is empty.

Here are the parameter details. Write the domain as \(D(\bar x,\bar p)\)
and the relation as \(\theta(\bar x,\bar y,\bar p)\), for a finite
parameter tuple \(\bar p\). Let \(B(\bar u)\) be the first-order
formula saying that \(\theta(-,-,\bar u)\) is reflexive, symmetric
and transitive on \(D(-,\bar u)\). Define on the whole product of
\(\bar x,\bar u\) blocks
\[
\begin{split}
 F((\bar x,\bar u),(\bar y,\bar v))\ \Longleftrightarrow\quad
 &\bar u=\bar v\ \land\bigl[
 B(\bar u)\land D(\bar x,\bar u)\land D(\bar y,\bar u)
                  \land\theta(\bar x,\bar y,\bar u)\\
 &\quad\lor\bigl(\neg(B(\bar u)\land D(\bar x,\bar u))
       \land\neg(B(\bar u)\land D(\bar y,\bar u))
       \land\bar x=\bar y\bigr)\bigr].
\end{split}
\]
Each fixed \(\bar u\)-fiber consists of its valid equivalence classes
and singleton outside tuples; an invalid fiber consists entirely of
singletons. Different fibers are never equivalent. Hence \(F\) is a
parameter-free equivalence relation in every source model. If \(B(\bar p)\)
holds, the map \([\bar a]_\theta\mapsto[(\bar a,\bar p)]_F\)
identifies the parameter-defined quotient with the definable subset obtained
by requiring the last representative block to equal \(\bar p\) and
the first to satisfy \(D\). This proves the claim without treating all
parameter choices as if they defined equivalence relations.

Finally pull back a parameter-free definable domain and equivalence relation
on a finite product of imaginary sorts to their finite home blocks. By the
truth lemma the pulled-back relation is an equivalence on the pulled-back
domain and is invariant under all initial quotient equalities. Extend it by
singleton classes outside that domain as above. Its sort already belongs
to \(M^{\mathrm{eq}}\). Two imaginary tuples have the same new code
exactly when they are equivalent, since changing any of their home
representatives does not change that code. The induced quotient map and its
inverse relation are definable by existentially quantifying the finitely many
home blocks and their projection equations. It is a bijection to the
specified definable subset. Thus closure under quotients uses the expanded
sorts, not an unproved real-coding property of the original structure.
\end{proof}
\end{annexitem}

\begin{annexitem}{annex:p03-imaginaries-canonical-code}{thm:p03-imaginaries-canonical-code}{Family equivalence and extension of automorphisms}
\begin{proof}
Equality of the sets defined by two parameter tuples is an equivalence
relation, and its first-order formula is exactly \(E_\varphi\).
Every tuple in the class of \(\bar a\) defines the same set as
\(\bar a\); the displayed existential formula therefore defines exactly
\(X\), with both directions witnessed by that equality of sets and the
representative \(\bar a\) itself.

An automorphism \(\sigma\) of the home structure preserves each
parameter-free \(E\), so \([\bar a]_E\mapsto[\sigma\bar a]_E\)
is well defined and bijective. It preserves projection maps and every
declared invariant relation by their home definitions, so is an automorphism
of the expansion. Surjectivity of every projection forces this action for
any extension, giving uniqueness. Now \(\sigma e_X=e_X\) is
\(E_\varphi(\sigma\bar a,\bar a)\), exactly the equality of
\(\varphi(M,\sigma\bar a)=\sigma X\) with \(X\). If the language
names parameters, the same calculation uses automorphisms fixing those names.
\end{proof}
\end{annexitem}

\begin{annexitem}{annex:p03-imaginaries-uniform-codes}{thm:p03-imaginaries-uniform-codes}{Definable image and stabilizer implications}
\begin{proof}
The image of \(f_E\) is defined by
\(\exists\bar x(D(\bar x)\land\chi_E(\bar x,\bar z))\).
Invariance on classes gives a map from \(D/E\) to this image; the
equivalence-of-equalities condition makes it injective, and its definition
makes it onto. Its graph in the imaginary expansion is
\[
 \exists\bar x(D(\bar x)\land q_{E^*}(\bar x)=e
                              \land\chi_E(\bar x,\bar z)).
\]
For each relevant \(e\) it has exactly one real tuple \(\bar z\);
for each tuple in the image it has exactly one \(e\). Thus the two
are interdefinable over the declared parameters. Uniqueness implies finiteness,
so this is also a weak code, including the zero-tuple case.

If \(s\) is constant on classes, \(E(x,y)\) implies \(s(x)=s(y)\).
Conversely \(s(x)=s(y)\), together with \(E(x,s(x))\) and
\(E(y,s(y))\), gives \(E(x,y)\) by symmetry and transitivity.
This proves the selector claim and shows why an arbitrary non-invariant
Skolem witness is insufficient.

A unique solution of a formula over \(A\) is fixed by every automorphism
fixing \(A\), and a finite solution set is permuted by those automorphisms.
Apply the uniqueness observation to the two defining formulas of an
interdefinable code: fixing either code forces fixing the other, so their
stabilizers over \(A\) agree. Apply the finite-set observation to a
formula witnessing algebraicity: its orbit is contained in that finite set.
For a finite tuple the product of its finite coordinate solution sets is
finite and definable. No converse from arbitrary automorphism-fixed points
to definable closure in a nonsaturated structure is being assumed.
\end{proof}
\end{annexitem}

\begin{annexitem}{annex:p03-imaginaries-unordered-pairs}{thm:p03-imaginaries-unordered-pairs}{Finite ambiguity versus a definable ordering}
\begin{proof}
The quotient projection makes \(e\) definable from \((a,b)\).
The formula \(q_{E^*}(x,y)=e\land x\ne y\) has exactly the two
tuple solutions \((a,b),(b,a)\). Hence both coordinates and the tuple
are algebraic over \(e\), giving a weak code.

Every permutation of a pure set is an automorphism. Permutations preserving
\(\{a,b\}\) include the transposition of \(a,b\), so neither is
fixed by the stabilizer of \(e\). Every other element can be moved to
another element outside that pair while fixing the pair. There is thus no
real element fixed by all automorphisms fixing \(e\). A nonempty real
tuple definable from \(e\) would have every coordinate fixed, a contradiction.
An empty tuple cannot code \(e\): a permutation moving the pair moves
\(e\), whereas an imaginary definable over the empty set is fixed by all
automorphisms. This proves failure of definable real coding.

In the ordered structure define the output as \((x,y)\) if \(x<y\)
and as \((y,x)\) if \(y<x\). This is a definable function on distinct
pairs, is constant on their unordered classes, and two outputs agree exactly
when the unordered pairs agree. The uniform coding criterion applies.
\end{proof}
\end{annexitem}

\begin{annexitem}{annex:p03-imaginaries-infinite-classes}{thm:p03-imaginaries-infinite-classes}{Infinite stabilizer orbits rule out algebraic codes}
\begin{proof}
For any real point \((i,j)\), arbitrarily many permutations within its
infinite \(R\)-class move it to distinct points, fixing every class
setwise and therefore fixing \(e\). It has an infinite orbit under the
stabilizer of \(e\). The algebraic-orbit implication in the preceding
proof shows that no real point is algebraic over \(e\). Consequently no
nonempty real tuple is algebraic over \(e\). The empty tuple is not a
weak code either, since a permutation exchanging the zeroth and first classes
is an automorphism and moves \(e\); thus \(e\) is not definable over
the empty set. This exhausts all finite real tuples.
\end{proof}
\end{annexitem}

\begin{annexitem}{annex:p03-imaginaries-quotient-group}{thm:p03-imaginaries-quotient-group}{Normality, descent, and a nonnormal counterexample}
\begin{proof}
Reflexivity, symmetry and transitivity of \(E\) use, respectively,
\(1\in N\), closure under inverse, and closure under product. If
\(x'=xn\), \(y'=ym\), with \(n,m\in N\), then
\[
 x'y'=xnym=xy(y^{-1}ny)m,
\]
and \((y^{-1}ny)m\in N\) by normality. Products of equivalent
inputs are therefore equivalent. For inverses,
\(x'^{-1}=n^{-1}x^{-1}=x^{-1}(xn^{-1}x^{-1})\), also equivalent
to \(x^{-1}\). Changing the output representative preserves the displayed
graphs by equivalence; each graph has an output and one output class.
Thus these are the required quotient data. Associativity, identity and inverse
laws follow by applying the quotient map to the corresponding group laws.
The class of \(x\) is precisely the coset \(xN\).

For necessity in the stated construction, take \(G=S_3\) and
\(N=\{1,(12)\}\), a definable subgroup with its two elements named.
Then \(1\) is equivalent to \((12)\), but multiplication on the right
by \((13)\) gives \((13)\) and \((12)(13)\), which are not
equivalent: their quotient difference is
\((13)(12)(13)=(23)\notin N\). Thus the multiplication graph is
not invariant. This counterexample retains a genuine equivalence relation
while rejecting the missing normality hypothesis.
\end{proof}
\end{annexitem}

\begin{annexitem}{annex:p03-imaginaries-fraction-field}{thm:p03-imaginaries-fraction-field}{Cross multiplication, field laws, and coding in a field}
\begin{proof}
Reflexivity and symmetry of \(E\) are immediate. If \(ad=cb\)
and \(cf=ed\), multiply the first equality by \(f\) and the second
by \(b\) to obtain \(adf=bed\). Since \(d\ne0\), cancellation
in the integral domain gives \(af=eb\), proving transitivity.
Products of nonzero denominators are nonzero, so all displayed outputs lie
in \(D\).

If \(ab'=a'b\) and \(cd'=c'd\), then
\[
 (ad+bc)b'd'=(a'd'+b'c')bd,\qquad
                       acb'd'=a'c'bd.
\]
These equations prove invariance of addition and multiplication. Negation
preserves the same cross-multiplication identity. Also \([a,b]=0\)
iff \(a=0\); if \([a,b]=[c,d]\) and \(a=0\), cancellation of
\(b\ne0\) forces \(c=0\). For nonzero classes the inverse is defined,
and \(ad=cb\) gives \(bc=da\), the equality of the inverse pairs.
The graphs are consequently invariant, total on their asserted domains and
single-valued on classes; inverse may be totalized at zero by defining its
value to be zero if the field signature requires a total inverse symbol.

The additive identity and inverse follow from
\([a,b]+[0,1]=[a,b]\) and
\([a,b]+[-a,b]=[0,b^2]=[0,1]\).
Addition is commutative. The two associations of three summands have a common
denominator \(bdf\) and numerator \(adf+bcf+bde\), proving
associativity. Multiplication is commutative and associative by its numerator
and denominator products, and \([1,1]\) is its identity. Distributivity
uses
\[
 \frac a b\left(\frac c d+\frac e f\right)
       =\frac{a(cf+de)}{bdf},\qquad
 \frac{ac}{bd}+\frac{ae}{bf}
       =\frac{ab(cf+de)}{b^2df},
\]
whose pairs are equivalent by cross multiplication. Here fraction notation
abbreviates the classes already constructed; it does not assume a field in
advance. For \(a\ne0\), \([a,b][b,a]=[ab,ba]=[1,1]\).
Finally \([0,1]\ne[1,1]\) because \(R\) is nonzero. These are all
commutative-field axioms. The map \(a\mapsto[a,1]\) preserves both
ring operations and constants by their formulas, and is injective because
\((a,1)E(c,1)\) iff \(a=c\).

When \(R\) is a field, \(bz=a\) has a unique solution for \(b\ne0\).
Two pairs have the same solution exactly when \(ad=cb\): substitute
\(a=bz,c=dz\) for one direction, and cancel the nonzero \(bd\)
for the other. Thus this fixed formula is a uniform real code. The existence
of this solution was used only in the field case, not to construct the
quotient over an arbitrary integral domain.
\end{proof}
\end{annexitem}
